%  LaTeX support: latex@mdpi.com 
%  For support, please attach all files needed for compiling as well as the log file, and specify your operating system, LaTeX version, and LaTeX editor.

%=================================================================
\documentclass[journal,article,submit,pdftex,moreauthors]{Definitions/mdpi} 
%\documentclass[preprints,article,submit,pdftex,moreauthors]{Definitions/mdpi} 
% For posting an early version of this manuscript as a preprint, you may use "preprints" as the journal. Changing "submit" to "accept" before posting will remove line numbers.

%--------------------
% Class Options:
%--------------------
%----------
% journal
%----------
% Choose between the following MDPI journals:
% accountaudit, acoustics, actuators, addictions, adhesives, adolescents, aerobiology, aerospace, agriculture, agriengineering, agrochemicals, agronomy, ai, aichem, aieng, aimater, aimed, aipa, air, aisens, algorithms, allergies, alloys, amh, analog, analytica, analytics, anatomia, anesthres, animals, antibiotics, antibodies, antioxidants, applbiosci, appliedchem, appliedmath, appliedphys, applmech, applmicrobiol, applnano, applsci, aquacj, architecture, arm, arthropoda, asc, asi, astronautics, astronomy, atmosphere, atoms, audiolres, automation, axioms, bacteria, batteries, bdcc, beverages, biochem, bioengineering, biologics, biology, biomass, biomechanics, biomed, biomedicines, biomedinformatics, biomimetics, biomolecules, biophysica, bioresourbioprod, biosensors, biosphere, biotech, birds, blockchains, bloods, blsf, brainsci, breath, buildings, cancers, carbon, cardio, cardiogenetics, cardiovascmed, catalysts, cells, ceramics, challenges, chemengineering, chemistry, chemosensors, chemproc, children, chips, cimb, civileng, cleantechnol, climate, clinbioenerg, clinpract, clockssleep, cmd, cmtr, coasts, coatings, colloids, colorants, commodities, complexities, complications, compounds, computation, computers, condensedmatter, conservation, constrmater, cosmetics, covid, crops, cryo, cryptography, crystals, csmf, ctn, culture, curroncol, cyber, dairy, data, ddc, dentistry, dermato, dermatopathology, designs, devices, dhi, diabetology, diagnostics, dietetics, digital, disabilities, diseases, diversity, dna, drones, dynamics, earth, ebj, ecm, ecologies, edm, eesp, electricity, electrochem, electronicmat, electronics, encyclopedia, endocrines, energies, eng, engproc, entomology, entropy, environments, environremediat, epidemiologia, epigenomes, esa, est, fermentation, fibers, fintech, fire, fishes, fluids, foods, forecasting, forensicsci, forests, fossstud, foundations, fractalfract, fuels, future, futureinternet, futurepharmacol, futurephys, futuretransp, galaxies, gases, gastroent, gastrointestdisord, gastronomy, gels, genes, geographies, geohazards, geomatics, geometry, geosciences, geotechnics, geriatrics, germs, glacies, grasses, green, greenhealth, gucdd, hardware, hazardousmatters, healthcare, hearts, hemato, hematolrep, hep, heritage, higheredu, highthroughput, horticulturae, hospitals, hydrobiology, hydrogen, hydrology, hydropower, hygiene, idr, iic, ijem, ijerph, ijgi, ijmd, ijms, ijns, ijom, ijpb, ijt, ijtm, ijtpp, ime, immuno, informatics, information, infrastructures, inorganics, insects, instruments, inventions, iot, j, jaestheticmed, jal, jcdd, jcm, jcp, jcrm, jcs, jcto, jdad, jdb, jdream, jemr, jeta, jfb, jfmk, jgbg, jgg, jimaging, jlpea, jmahp, jmmp, jmms, jmp, jmse, jne, jnt, jof, joi, joitmc, joma, jop, joptm, jor, jox, jpbi, jphytomed, jpm, jsan, jtaer, jvd, jzbg, kidneydial, kinasesphosphatases, knowledge, labmed, laboratories, lae, land, life, lights, limnolrev, lipidology, liquids, livers, logics, logistics, lubricants, lymphatics, machines, macromol, magnetism, magnetochemistry, make, marinedrugs, materials, materproc, mathematics, mca, measurements, medicina, medicines, medsci, membranes, merits, metabolites, metals, meteorology, methane, metrics, metrology, micro, microarrays, microbiolres, microelectronics, micromachines, microorganisms, microplastics, microwave, minerals, mining, mmphys, modelling, molbank, molecules, mps, msf, mti, multimedia, muscles, nanoenergyadv, nanomanufacturing, nanomaterials, ncrna, ndt, network, neuroglia, neuroimaging, neurolint, neurosci, nitrogen, notspecified, nursrep, nutraceuticals, nutrients, obesities, occuphealth, oceans, ohbm, onco, optics, oral, organics, organoids, osteology, oxygen, pandemics, parasites, parasitologia, particles, pathogens, pathophysiology, pediatrrep, pets, pharmaceuticals, pharmaceutics, pharmacoepidemiology, pharmacy, philosophies, photochem, photonics, phycology, physchem, physics, physiologia, plants, plasma, platforms, pollutants, polymers, polysaccharides, populations, poultry, powders, precipitation, precisoncol, preprints, proceedings, processes, prosthesis, proteomes, psf, psychiatryint, psychoactives, purification, quantumrep, quaternary, qubs, radiation, rdt, reactions, realestate, receptors, recycling, regeneration, remotesensing, reports, reprodmed, resources, rheumato, rjpm, robotics, rsee, ruminants, safety, sci, scipharm, sclerosis, seeds, sensors, separations, sexes, shi, signals, sinusitis, siuj, skins, smartcities, sna, societies, software, soilsystems, solar, solids, spectroscj, sports, standards, stats, std, stratsediment, stresses, surfaces, surgeries, suschem, sustainability, symmetry, synbio, systems, tae, targets, taxonomy, technologies, telecom, test, textiles, thalassrep, therapeutics, thermo, timespace, tomography, toxics, toxins, tph, transplantology, transportation, traumacare, traumas, tri, tropicalmed, universe, urbansci, uro, vaccines, vehicles, venereology, vetsci, vibration, virtualworlds, viruses, vision, waste, water, welding, wem, wevj, wild, wind, women, world, zoonoticdis

%---------
% article
%---------
% The default type of manuscript is "article", but can be replaced by: 
% abstract, addendum, article, benchmark, book, bookreview, briefcommunication, briefreport, casereport, changes, clinicopathologicalchallenge, comment, commentary, communication, conceptpaper, conferenceproceedings, correction, conferencereport, creative, datadescriptor, discussion, entry, expressionofconcern, extendedabstract, editorial, essay, erratum, fieldguide, hypothesis, interestingimages, letter, meetingreport, monograph, newbookreceived, obituary, opinion, proceedingpaper, projectreport, reply, retraction, review, perspective, protocol, shortnote, studyprotocol, supfile, systematicreview, technicalnote, viewpoint, guidelines, registeredreport, tutorial,  giantsinurology, urologyaroundtheworld
% supfile = supplementary materials

%----------
% submit
%----------
% The class option "submit" will be changed to "accept" by the Editorial Office when the paper is accepted. This will only make changes to the frontpage (e.g., the logo of the journal will get visible), the headings, and the copyright information. Also, line numbering will be removed. Journal info and pagination for accepted papers will also be assigned by the Editorial Office.

%------------------
% moreauthors
%------------------
% If there is only one author the class option oneauthor should be used. Otherwise use the class option moreauthors.

%---------
% pdftex
%---------
% The option pdftex is for use with pdfLaTeX. Remove "pdftex" for (1) compiling with LaTeX & dvi2pdf (if eps figures are used) or for (2) compiling with XeLaTeX.

%=================================================================
% MDPI internal commands - do not modify
\firstpage{1} 
\makeatletter 
\setcounter{page}{\@firstpage} 
\makeatother
\pubvolume{1}
\issuenum{1}
\articlenumber{0}
\pubyear{2026}
\copyrightyear{2026}
%\externaleditor{Firstname Lastname} % More than 1 editor, please add `` and '' before the last editor name
\datereceived{ } 
\daterevised{ } % Comment out if no revised date
\dateaccepted{ } 
\datepublished{ } 
%\datecorrected{} % For corrected papers: "Corrected: XXX" date in the original paper.
%\dateretracted{} % For retracted papers: "Retracted: XXX" date in the original paper.
%\doinum{} % Used for some special journals, like molbank
%\pdfoutput=1 % Uncommented for upload to arXiv.org
%\CorrStatement{yes}  % For updates
%\longauthorlist{yes} % For many authors that exceed the left citation part
%\IsAssociation{yes} % For association journals

%=================================================================
% Add packages and commands here. The following packages are loaded in our class file: fontenc, inputenc, calc, indentfirst, fancyhdr, graphicx, epstopdf, lastpage, ifthen, float, amsmath, amssymb, lineno, setspace, enumitem, mathpazo, booktabs, titlesec, etoolbox, tabto, xcolor, colortbl, soul, multirow, microtype, tikz, totcount, changepage, attrib, upgreek, array, tabularx, pbox, ragged2e, tocloft, marginnote, marginfix, enotez, amsthm, natbib, hyperref, cleveref, scrextend, url, geometry, newfloat, caption, draftwatermark, seqsplit
% cleveref: load \crefname definitions after \begin{document}

%=================================================================
% Please use the following mathematics environments: Theorem, Lemma, Corollary, Proposition, Characterization, Property, Problem, Example, ExamplesandDefinitions, Hypothesis, Remark, Definition, Notation, Assumption
%% For proofs, please use the proof environment (the amsthm package is loaded by the MDPI class).

%=================================================================
% Full title of the paper (Capitalized)
\Title{Title}

% Author Orchid ID: enter ID or remove command
\newcommand{\orcidauthorA}{0000-0000-0000-000X} % Add \orcidA{} behind the author's name
%\newcommand{\orcidauthorB}{0000-0000-0000-000X} % Add \orcidB{} behind the author's name

% Authors, for the paper (add full first names)
\Author{Catalina Valles González $^{1}$\orcidA{}*}

%\longauthorlist{yes}

% MDPI internal command: Authors, for metadata in PDF
\AuthorNames{Firstname Lastname, Firstname Lastname and Firstname Lastname}

% Affiliations / Addresses (Add [1] after \address if there is only one affiliation.)
\address{%
$^{1}$ \quad Affiliation 1; e-mail@e-mail.com\\
$^{2}$ \quad Affiliation 2; e-mail@e-mail.com}

% Contact information of the corresponding author
\corres{Correspondence: e-mail@e-mail.com; Tel.: (optional; include country code; if there are multiple corresponding authors, add author initials) +xx-xxxx-xxx-xxxx (F.L.)}

% Current address and/or shared authorship
%\firstnote{Current address: Affiliation.}  
% Current address should not be the same as any items in the Affiliation section.

%\secondnote{These authors contributed equally to this work.}
% The commands \thirdnote{} till \eighthnote{} are available for further notes.

%\simplesumm{} % Simple summary

%\conference{} % An extended version of a conference paper

% Abstract (Do not insert blank lines, i.e. \\) 
\abstract{Demand forecasting plays a critical role in operations management, particularly in environments where production and inventory decisions depend on forward-looking estimates. However, many organizations still rely on simple averages and manual calculations, limiting forecast accuracy and systematic error evaluation. This study addresses this gap by evaluating nine forecasting methods—simple average, moving average, weighted average, simple exponential smoothing, double exponential smoothing, Holt–Winters, linear regression, ARIMA, and SARIMA—using error metrics within a walk-forward validation scheme to identify the most appropriate technique for different demand series. Based on this comparative analysis, a computational tool was developed in Python using the Dash library to automate data loading, method implementation, metric calculation, and interactive visualization of results. Additionally, the computational performance of the tool was assessed by analyzing processing times under different time series sizes. The results show that integrating multiple forecasting methods within a unified and automated environment enables systematic comparison, improves transparency, and supports more informed production and inventory decisions. The study concludes that computational tools that combine standardized evaluation metrics and interactive visualization can strengthen demand planning processes and reduce reliance on empirical forecasting practices.}

% Keywords
\keyword{Demand forecasting; time series; statistical models; production planning; Python Dash; data analysis; model optimization.} 

% The fields PACS, MSC, and JEL may be left empty or commented out if not applicable
%\PACS{J0101}
%\MSC{}
%\JEL{}

%%%%%%%%%%%%%%%%%%%%%%%%%%%%%%%%%%%%%%%%%%
% Only for the journal Diversity
%\LSID{\url{http://}}

%%%%%%%%%%%%%%%%%%%%%%%%%%%%%%%%%%%%%%%%%%
% Only for the journal Applied Sciences
%\featuredapplication{Authors are encouraged to provide a concise description of the specific application or a potential application of the work. This section is not mandatory.}
%%%%%%%%%%%%%%%%%%%%%%%%%%%%%%%%%%%%%%%%%%

%%%%%%%%%%%%%%%%%%%%%%%%%%%%%%%%%%%%%%%%%%
% Only for the journal Data
%\dataset{DOI number or link to the deposited data set if the data set is published separately. If the data set shall be published as a supplement to this paper, this field will be filled by the journal editors. In this case, please submit the data set as a supplement.}
%\datasetlicense{License under which the data set is made available (CC0, CC-BY, CC-BY-SA, CC-BY-NC, etc.)}

%%%%%%%%%%%%%%%%%%%%%%%%%%%%%%%%%%%%%%%%%%
% Only for the journal BioTech, Fishes, Neuroimaging and Toxins
%\keycontribution{The breakthroughs or highlights of the manuscript. Authors can write one or two sentences to describe the most important part of the paper.}

%%%%%%%%%%%%%%%%%%%%%%%%%%%%%%%%%%%%%%%%%%
% Only for the journal Encyclopedia
%\encyclopediadef{For entry manuscripts only: please provide a brief overview of the entry title instead of an abstract.}

%%%%%%%%%%%%%%%%%%%%%%%%%%%%%%%%%%%%%%%%%%
% Different journals have different requirements. Please check the specific journal guidelines in the "Instructions for Authors" on the journal's official website.
%\addhighlights{yes}
%\renewcommand{\addhighlights}{%
%
%\noindent The goal is to increase the discoverability and readability of the article via search engines and other scholars. Highlights should not be a copy of the abstract, but a simple text allowing the reader to quickly and simplified find out what the article is about and what can be cited from it. Each of these parts should be devoted up to 2~bullet points.\vspace{3pt}\\
%\textbf{What are the main findings?}
% \begin{itemize}[labelsep=2.5mm,topsep=-3pt]
% \item First bullet.
% \item Second bullet.
% \end{itemize}\vspace{3pt}
%\textbf{What are the implications of the main findings?}
% \begin{itemize}[labelsep=2.5mm,topsep=-3pt]
% \item First bullet.
% \item Second bullet.
% \end{itemize}
%}

%%%%%%%%%%%%%%%%%%%%%%%%%%%%%%%%%%%%%%%%%%
\begin{document}

%%%%%%%%%%%%%%%%%%%%%%%%%%%%%%%%%%%%%%%%%%
\setcounter{section}{-1} %% Remove this when starting to work on the template.
\section{How to Use this Template}

The template details the sections that can be used in a manuscript. Note that the order and names of article sections may differ from the requirements of the journal (e.g., the positioning of the Materials and Methods section). Please check the instructions on the authors' page of the journal to verify the correct order and names. For any questions, please contact the editorial office of the journal or support@mdpi.com. For LaTeX-related questions please contact latex@mdpi.com.%\endnote{This is an endnote.} % To use endnotes, please un-comment \printendnotes below (before References). Only journal Laws uses \footnote.

% The order of the section titles is different for some journals. Please refer to the "Instructions for Authors” on the journal homepage.

\section{Introduction}
"((((ELIMINAR
The introduction should briefly place the study in a broad context and highlight why it is important. It should define the purpose of the work and its significance. The current state of the research field should be reviewed carefully and key publications cited. Please highlight controversial and diverging hypotheses when necessary. Finally, briefly mention the main aim of the work and highlight the principal conclusions. As far as possible, please keep the introduction comprehensible to scientists outside your particular field of research. Citing a journal paper  \cite{ref-journal}.  Now citing a book reference  \cite{ref-book1,ref-book2} or other reference types  \cite{ref-unpublish,ref-url}. Please use the command  \cite{ref-proceeding,ref-thesis} for the following MDPI journals, which use author--date citation: Administrative Sciences, Arts, Behavioral Sciences, Businesses, Econometrics, Economies, Education Sciences, European Journal of Investigation in Health, Psychology and Education, Games, Genealogy, Histories, Humanities, Humans, IJFS, Journal of Intelligence, Journalism and Media, JRFM, Languages, Laws, Literature, Psychology International, Publications, Religions, Risks, Social Sciences, Tourism and Hospitality, Youth. 
))))"

Demand planning plays a central role in operations management, particularly in environments where production, procurement, and inventory decisions depend on anticipatory information. In increasingly dynamic markets characterized by volatility, global competition, and evolving consumption patterns, forecasting future demand becomes both complex and indispensable (Mentzer et al., 2001; Chopra and Meindl, 2021). Accurate forecasting reduces uncertainty, improves resource allocation, and mitigates operational imbalances. However, demand rarely follows perfectly stable patterns. Seasonality, structural variability, irregular consumption behavior, and external shocks limit predictive precision and challenge traditional forecasting methods (Makridakis et al., 1998; Hyndman and Athanasopoulos, 2018).

The importance of forecasting in logistics and inventory systems has been extensively documented. Forecasts directly influence storage decisions, transportation planning, procurement policies, and safety stock levels (Ballou, 2004). Persistent forecast errors may generate overstocks, stockouts, and service-level instability (Silver, Pyke and Peterson, 1998). Empirical studies show that improving forecast accuracy can significantly reduce total inventory costs while increasing product availability, highlighting the strong relationship between forecast accuracy and operational performance (Syntetos et al., 2009; Fildes et al., 2019). Recent research further indicates that enhanced forecasting performance may translate into measurable cost reductions and improved supply chain responsiveness (Sattar et al., 2025). 

From a methodological perspective, classical time-series models such as moving averages, exponential smoothing, regression models, and ARIMA-based approaches remain widely used due to their interpretability and ability to capture trend and seasonality components (Makridakis et al., 1998; Nahmias and Olsen, 2015; Hyndman and Athanasopoulos, 2018).

Nonetheless, in environments characterized by intermittent, highly volatile, or structurally unstable demand, these traditional models may underperform. This has encouraged the integration of advanced computational tools, including machine learning algorithms and data-driven forecasting systems (Carbonneau et al., 2008; Luo, 2024).
Despite methodological advances, organizational practice often lags behind academic development. Many firms continue to rely on manual spreadsheets and empirical adjustments, limiting systematic error evaluation and reducing transparency in forecasting processes (Madariaga Fernández, 2020). Kerkkänen, Korpela, and Huiskonen (2009) report that a substantial proportion of companies fail to properly measure forecast errors, weakening the integration between forecasting and operational planning. This gap between theoretical capability and practical implementation remains a persistent challenge in supply chain management.

At the same time, the literature reflects diverging perspectives regarding the superiority of complex models over simpler statistical techniques. Large-scale forecasting competitions (e.g., the M-competitions) have shown that, in many contexts, relatively simple statistical methods perform comparably to or better than more sophisticated approaches (Makridakis et al., 2018). Conversely, other studies argue that machine learning and hybrid models outperform traditional techniques when large datasets and nonlinear patterns are present (Carbonneau et al., 2008). This ongoing debate suggests that forecasting performance is highly context-dependent and reinforces the need for systematic model comparison under homogeneous evaluation criteria.

Recent research also emphasizes that computational tools enhance decision-making quality when uncertainty is explicitly visualized and interpreted. Visual analytics approaches integrating automated statistical modeling with interactive dashboards improve user understanding and facilitate managerial adoption (Keim et al., 2010; Maack et al., 2025). In particular, uncertainty-aware visualization frameworks allow decision-makers to assess forecast variability, compare competing models, and better understand error behavior. These developments support the design of accessible computational tools capable of integrating statistical rigor with intuitive interfaces.

In this context, this study addresses the following research question: Which forecasting methods provide superior performance under varying demand structures when evaluated through consistent statistical metrics, and how can this evaluation process be automated through an interactive computational tool? To answer this question, nine forecasting methods are implemented and compared: simple average, moving average, weighted moving average, simple exponential smoothing, double exponential smoothing (Holt), Holt–Winters, linear regression, ARIMA, and SARIMA. Each method is evaluated using standardized performance metrics, including Mean Error (ME), Mean Absolute Deviation (MAD), Mean Squared Error (MSE), and Mean Absolute Percentage Error (MAPE), ensuring objective comparability.

The primary contribution of this study lies in the development of a generalized computational framework that integrates multiple classical forecasting methods, automated error evaluation, and interactive visualization in a single platform. The tool is developed in Python using the Dash library, enabling users without advanced statistical training to upload historical series, compare models under identical conditions, and interpret uncertainty through dynamic visualizations.
In addition to methodological comparison, computational performance is analyzed by measuring processing time as a function of series length, providing insight into scalability and practical feasibility. The principal findings indicate that forecasting accuracy varies significantly according to demand structure, that no single model universally dominates across scenarios, and that automated comparative evaluation substantially improves transparency and decision support. These results reinforce the importance of structured forecasting processes over reliance on isolated or manual techniques.

By combining statistical modeling, systematic evaluation, and interactive visualization, this study contributes to bridging the gap between forecasting theory and operational practice. The proposed framework provides a replicable and scalable approach for organizations seeking to improve demand planning robustness while maintaining methodological transparency.

Finally, this paper is organized as follows. Section 2 describes the research methodology, including model implementation, evaluation metrics, and computational design. Section 3 presents the empirical and computational results. Section 4 discusses findings in relation to the existing literature. Section 5 concludes with implications, limitations, and directions for future research.

%%%%%%%%%%%%%%%%%%%%%%%%%%%%%%%%%%%%%%%%%%
\section{Materials and Methods}
((( ELIMINAR
Materials and Methods should be described with sufficient details to allow others to replicate and build on published results. Please note that publication of your manuscript implicates that you must make all materials, data, computer code, and protocols associated with the publication available to readers. Please disclose at the submission stage any restrictions on the availability of materials or information. New methods and protocols should be described in detail while well-established methods can be briefly described and appropriately cited.

Research manuscripts reporting large datasets that are deposited in a publicly avail-able database should specify where the data have been deposited and provide the relevant accession numbers. If the accession numbers have not yet been obtained at the time of submission, please state that they will be provided during review. They must be provided prior to publication.

Interventionary studies involving animals or humans, and other studies require ethical approval must list the authority that provided approval and the corresponding ethical approval code.

In this section, where applicable, authors are required to disclose details of how gen-erative artificial intelligence (GenAI) has been used in this paper (e.g., to generate text, data, or graphics, or to assist in study design, data collection, analysis, or interpretation). The use of GenAI for superficial text editing (e.g., grammar, spelling, punctuation, and formatting) does not need to be declared.

\begin{quote}
This is an example of a quote.
\end{quote}

)))

\subsection{Research Design}

This study follows a quantitative and applied research design aimed at mathematically analyzing demand time series and developing a generalized computational forecasting framework suitable for real-world operational environments. The research is non-experimental, as historical demand records are analyzed without manipulation of independent variables. It is longitudinal in nature, given that monthly time-series data are evaluated over multiple periods to capture structural dynamics.
The methodological approach is descriptive and propositional. First, structural components of the time series (trend, seasonality, stationarity) are identified. Second, a computational tool is developed to automate model comparison, parameter optimization, and forecast generation. Similar applied quantitative approaches are widely adopted in demand planning and inventory forecasting research (Petropoulos et al., 2022; Makridakis et al., 2022).

\subsection{Data and Experimental Material}

The experimental material consists of monthly demand time series structured with three variables: year, month, and demand. Although empirical validation was conducted using industrial sales data, the tool is fully generalized and accepts any dataset that follows this structure.
Input files must be provided in .xlsx format with the exact column names: year, month and demand.
During data loading, the system:Converts textual month labels into numerical format,
constructs a monthly time index, validates column presence and data coherence, checks numeric consistency of demand values. If validation fails, the system returns an error message, preventing further processing. This ensures data integrity and replicability.

\subsection{Series Classification Procedure}

Before model comparison, the tool performs structural classification along three dimensions:
\begin{itemize}
\item Trend Detection:
A simple linear regression of demand against time is estimated. Statistical significance of the slope (p < 0.05) indicates deterministic trend presence. This approach aligns with recent forecasting structure-detection practices (Talagala et al., 2021).
\item Seasonality Detection:
Autocorrelation at lag 12 is computed. Absolute autocorrelation above 0.30 indicates annual seasonality, consistent with time-series diagnostics standards (Wei, 2006; Montgomery et al., 2015).
\item Stationarity:
The Augmented Dickey–Fuller test (ADF) is applied. A p-value < 0.05 indicates stationarity.
The classification results are displayed to users but do not automatically transform the models, preserving methodological transparency.
\end{itemize}

\subsection{Forecasting Methods Included}

Nine forecasting methods are implemented:
\begin{enumerate}
\item Simple Average
\item Moving Average
\item Weighted Moving Average
\item Simple Exponential Smoothing
\item Holt’s Linear Trend
\item Holt–Winters
\item Linear Regression
\item ARIMA
\item SARIMA
\end{enumerate}
Classical models follow formulations described in Nahmias and Olsen (2015) and Montgomery et al. (2015).
The inclusion of both simple statistical models and ARIMA-based approaches is consistent with findings from large-scale forecasting competitions, demonstrating that no single method dominates across all demand structures (Makridakis et al., 2022).
Recent evidence shows that systematic comparison across heterogeneous models improves robustness under demand uncertainty (Petropoulos et al., 2022).

\subsection {Validation Strategy}

Model evaluation follows a walk-forward one-step-ahead rolling origin validation scheme, widely recommended for dependent time-series data (Bergmeir et al., 2020).
For each time step t:
\begin{itemize}
\item The model is trained using data up to \begin{equation} t-1 \end{equation}
\item A one-step-ahead forecast is generated
\item The process repeats recursively across the series
\end{itemize}
This procedure prevents look-ahead bias and yields realistic performance estimates, as recommended in contemporary forecasting evaluation research (Makridakis et al., 2022).

\subsection {Error Metrics}
The following metrics are computed:
\begin{itemize}
\item Mean Absolute Deviation (MAD)
\item Mean Squared Error (MSE)
\item Mean Absolute Percentage Error (MAPE)
\item Mean Error (ME)
\end{itemize}
MAPE is used for final ranking due to interpretability and comparability across scales (Makridakis et al., 2022).

\subsection {Hyperparameter Optimization}
After initial ranking, the best-performing model undergoes parameter optimization through grid search:
\begin{itemize}
\item SES: alpha sweep
\item Holt: (alpha, beta) grid
\item Holt–Winters: (alpha, beta, gamma) two-stage grid
\item ARIMA/SARIMA: compact (p,d,q) and seasonal (P,D,Q) combinations
Each parameter combination is re-evaluated using walk-forward validation.
Search ranges are computationally bounded to prevent exponential growth in runtime, consistent with scalable forecasting system design principles (Benidis et al., 2022).
\end{itemize}

\subsection {Computational Implementation}

The tool was developed in Python using: pandas, numpy, statsmodels, plotly, dash, and dash bootstrap.
The architecture is modular:
\begin{itemize}
\item Data upload and validation
\item Series classification and model evaluation
\item Hyperparameter optimization
\item Final forecasting and export
\end{itemize}

The general structure is shown in Table 1, which summarizes what the user can do in each module.
\begin{table}[H]
\caption{Stages and Components of the Tool}
\centering
\begin{tabularx}{\textwidth}{lX}
\toprule
\textbf{Stage} & \textbf{Components} \\
\midrule
Data loading and validation & Upload the file in .xlsx format (performed by the user) and display the uploaded data graphically and in tabular format. \\
Method classification and evaluation & Perform automatic classification and evaluate forecasting methods, selecting the model with the lowest error metrics and displaying a graphical comparison of the top three models with the lowest errors. \\
Parameter optimization of the winning method & Optimize parameters to further reduce error metrics. \\
Forecasting with the selected method & Graphically display three components: historical demand, simulated forecast over the historical demand, and future forecast. Additionally, provide a summary table of forecasted values and a button at the end to download the generated forecast. \\
\bottomrule
\end{tabularx}
\end{table}

\begin{figure}[H]
    \includegraphics[width=0.8\textwidth]{flowchart_tool.png}
    \caption{Methodological flow of the forecasting tool}
    \label{fig:demanda}
\end{figure}

\subsection{Data, Code, and Availability}

All computer code associated with this study will be made publicly available in a GitHub repository upon acceptance. The repository will include: Source code, example datasets, execution instructions, and reproducibility documentation.

\subsection{Use of Generative Artificial Intelligence (GenAI)}
Generative Artificial Intelligence tools were used to assist in structuring portions of the Python code based on predefined methodological specifications and supporting literature organization and preliminary synthesis using references explicitly selected and provided by the author.
All methodological decisions, model selection criteria, validation design, parameter ranges, interpretation of results, and final manuscript content were defined, verified, and validated exclusively by the author. No automated systems were used for data generation, statistical estimation, or independent analytical interpretation.

%%%%%%%%%%%%%%%%%%%%%%%%%%%%%%%%%%%%%%%%%%
\section{Results}

(((( ELIMINAR
This section may be divided by subheadings. It should provide a concise and precise description of the experimental results, their interpretation as well as the experimental conclusions that can be drawn.
\subsection{Subsection}
\subsubsection{Subsubsection}

Bulleted lists look like this:
\begin{itemize}
\item	First bullet;
\item	Second bullet;
\item	Third bullet.
\end{itemize}

Numbered lists can be added as follows:
\begin{enumerate}
\item	First item; 
\item	Second item;
\item	Third item.
\end{enumerate}

The text continues here.

\subsection{Figures, Tables and Schemes}

All figures and tables should be cited in the main text as Figure~\ref{fig1}, Table~\ref{tab1}, etc.

\begin{figure}[H]
%\isPreprints{\centering}{} % Only used for preprints

\includegraphics[width=4.0 cm]{Definitions/logo-mdpi}
\caption{This is a figure. Schemes follow the same formatting.\label{fig1}}
\end{figure}   
\unskip

\begin{table}[H] 
%\small % Change table font size
\caption{This is a table caption. Tables should be placed in the main text near to the first time they are~cited.\label{tab1}}
%\isPreprints{\centering}{} % Only used for preprints
\begin{tabularx}{\textwidth}{CCC}
\toprule
\textbf{Title 1}	& \textbf{Title 2}	& \textbf{Title 3}\\
\midrule
Entry 1		& Data			& Data\\
Entry 2		& Data			& Data \textsuperscript{1}\\
\bottomrule
\end{tabularx}

\noindent{\footnotesize{\textsuperscript{1} Tables may have a footer.}}
\end{table}

The text continues here (Figure~\ref{fig2} and Table~\ref{tab2}).

% Example of a figure that spans the whole page width and with subfigures. The same concept works for tables, too.
\begin{figure}[H]
%\isPreprints{}{% This command is only used for ``preprints''.
\begin{adjustwidth}{-\extralength}{0cm}
\centering
%} % If the paper is ``preprints'', please uncomment this parenthesis.
\subfloat[\centering]{
\includegraphics[width=7.0cm]{Definitions/logo-mdpi}}
%\hfill
\subfloat[\centering]{
\includegraphics[width=7.0cm]{Definitions/logo-mdpi}}\\
\subfloat[\centering]{
\includegraphics[width=7.0cm]{Definitions/logo-mdpi}}
%\hfill
\subfloat[\centering]{
\includegraphics[width=7.0cm]{Definitions/logo-mdpi}}
%\isPreprints{}{% This command is only used for ``preprints''.
\end{adjustwidth}
%} % If the paper is ``preprints'', please uncomment this parenthesis.
\caption{This is a wide figure. Schemes follow the same formatting. If there are multiple panels, they should be listed as: (\textbf{a}) Description of what is contained in the first panel. (\textbf{b}) Description of what is contained in the second panel. (\textbf{c}) Description of what is contained in the third panel. (\textbf{d}) Description of what is contained in the fourth panel. Figures should be placed in the main text near to the first time they are cited. A caption on a single line should be centered.\label{fig2}}
\end{figure} 

\begin{table}[H]
\caption{This is a wide table.\label{tab2}}
%\isPreprints{\centering}{% This command is only used for ``preprints''.
	\begin{adjustwidth}{-\extralength}{0cm}
%} % If the paper is ``preprints'', please uncomment this parenthesis.
%\isPreprints{\begin{tabularx}{\textwidth}{CCCC}}{% This command is only used for ``preprints''.
		\begin{tabularx}{\fulllength}{CCCC}
%} % If the paper is ``preprints'', please uncomment this parenthesis.
			\toprule
			\textbf{Title 1}	& \textbf{Title 2}	& \textbf{Title 3}     & \textbf{Title 4}\\
			\midrule
\multirow[m]{3}{*}{Entry 1 *}	& Data			& Data			& Data\\
			  	                   & Data			& Data			& Data\\
			             	      & Data			& Data			& Data\\
                   \midrule
\multirow[m]{3}{*}{Entry 2}    & Data			& Data			& Data\\
			  	                  & Data			& Data			& Data\\
			             	     & Data			& Data			& Data\\
			\bottomrule
		\end{tabularx}
%		\isPreprints{}{% This command is only used for ``preprints''.
	\end{adjustwidth}
%} % If the paper is ``preprints'', please uncomment this parenthesis.
	\noindent{\footnotesize{* Tables may have a footer.}}
\end{table}

%\begin{listing}[H]
%\caption{Title of the listing}
%\rule{\columnwidth}{1pt}
%\raggedright Text of the listing. In font size footnotesize, small, or normalsize. Preferred format: left aligned and single spaced. Preferred border format: top border line and bottom border line.
%\rule{\columnwidth}{1pt}
%\end{listing}

Text.

Text.

\subsection{Formatting of Mathematical Components}

This is the example 1 of equation:
\begin{linenomath}
\begin{equation}
a = 1,
\end{equation}
\end{linenomath}
the text following an equation need not be a new paragraph. Please punctuate equations as regular text.
%% If the documentclass option "submit" is chosen, please insert a blank line before and after any math environment (equation and eqnarray environments). This ensures correct linenumbering. The blank line should be removed when the documentclass option is changed to "accept" because the text following an equation should not be a new paragraph.

This is the example 2 of equation:
%\isPreprints{}{% This command is only used for ``preprints''.
\begin{adjustwidth}{-\extralength}{0cm}
%} % If the paper is ``preprints'', please uncomment this parenthesis.
\begin{equation}
a = b + c + d + e + f + g + h + i + j + k + l + m + n + o + p + q + r + s + t + u + v + w + x + y + z
\end{equation}
%\isPreprints{}{% This command is only used for ``preprints''.
\end{adjustwidth}
%} % If the paper is ``preprints'', please uncomment this parenthesis.

%% Example of a page in landscape format (with table and table footnote).
%\startlandscape
%\begin{table}[H] %% Table in wide page
%%\isPreprints{\centering}{} % This command is only used for ``preprints''.
%\caption{This is a very wide table.\label{tab3}}
%	\begin{tabularx}{\textwidth}{CCCC}
%		\toprule
%		\textbf{Title 1}	& \textbf{Title 2}	& \textbf{Title 3}	& \textbf{Title 4}\\
%		\midrule
%		Entry 1		& Data			& Data			& This cell has some longer content that runs over two lines.\\
%		Entry 2		& Data			& Data			& Data\textsuperscript{1}\\
%		\bottomrule
%	\end{tabularx}
%%\isPreprints{}{% This command is only used for ``preprints''.
%	\begin{adjustwidth}{+\extralength}{0cm}
%%} % If the paper is ``preprints'', please uncomment this parenthesis.
%		\noindent\footnotesize{\textsuperscript{1} This is a table footnote.}
%%\isPreprints{}{% This command is only used for ``preprints''.
%	\end{adjustwidth}
%%} % If the paper is ``preprints'', please uncomment this parenthesis.
%\end{table}
%\finishlandscape

Please punctuate equations as regular text. Theorem-type environments (including propositions, lemmas, corollaries etc.) can be formatted as follows:
%% Example of a theorem:
\begin{Theorem}
Example text of a theorem.
\end{Theorem}

The text continues here. Proofs must be formatted as follows:

%% Example of a proof:
\begin{proof}[Proof of Theorem 1]
Text of the proof. Note that the phrase ``of Theorem 1'' is optional if it is clear which theorem is being referred to.
\end{proof}
The text continues here.


ELIMINAR )))

\subsection{Tool Analysis}
To illustrate the functionality of the tool and evaluate its performance, the results obtained using the historical monthly demand of a dataset comprising 24 consecutive months of continuous information are presented. The analysis is structured according to the operational workflow of the tool: (1) Data loading and validation, (2) Method classification and evaluation, (3) Parameter optimization of the selected method, and (4) Forecasting using the selected method.

\subsubsection{Data Loading and Validation}

\begin{figure}[H]
    \includegraphics[width=0.8\textwidth]{flowchart_tool.png}
    \caption{Results of data loading and validation of the dataset}
    \label{fig:demanda}
\end{figure}

Once the series is uploaded into the tool, Figure NNNNNNN shows that the system displays a confirmation message indicating that the series is valid, as it contains 24 or more observations. Additionally, a plot of the uploaded series is presented, comprising 24 months of information between January 2023 and December 2024.
For the dataset analyzed in this case, the historical demand plot reveals variability over time, the presence of peaks in certain periods, and relatively abrupt changes between some months. Furthermore, the tool provides a tabular summary displaying the administered data.

\subsubsection{Method Classification and Evaluation}

\begin{figure}[H]
    \includegraphics[width=0.8\textwidth]{flowchart_tool.png}
    \caption{Classification results of the loaded data set}
    \label{fig:demanda}
\end{figure}

Once the series is loaded, Figure NNNNN presents the initial stage of Step 2 of the tool, where the automatic classification is displayed. In this case, the results indicate that the series exhibits a significant trend, does not show seasonal patterns, and is stationary according to the ADF test.

\begin{figure}[H]
    \includegraphics[width=0.8\textwidth]{flowchart_tool.png}
    \caption{Method evaluation for the loaded data set}
    \label{fig:demanda}
\end{figure}

Based on the previous classification, Figure 4 illustrates how the tool applies internal filtering rules to exclude forecasting methods that are structurally inconsistent with the identified characteristics of the series. Since the series does not exhibit seasonality, the Holt–Winters method and SARIMA models with seasonal components were excluded from evaluation. Conversely, because the series presents a statistically significant trend, models capable of capturing systematic level changes—such as Linear Regression, ARIMA, and Holt—were retained.
This selection strategy is consistent with recent forecasting research emphasizing that model performance is highly dependent on structural time-series features and that model selection should be aligned with these characteristics rather than applied uniformly across heterogeneous data (Montero-Manso et al., 2020; Makridakis et al., 2022). Feature-based model selection approaches demonstrate that incorporating structural information—such as trend or seasonality—improves forecast robustness and reduces the risk of structural mismatch.
In this case, Linear Regression achieved the best performance, presenting the lowest MAPE and greater error stability compared with Holt and ARIMA, indicating that deterministic trend modeling was sufficient to capture the dominant demand dynamics.

\begin{figure}[H]
    \includegraphics[width=0.8\textwidth]{flowchart_tool.png}
    \caption{Error comparison results for the loaded data set}
    \label{fig:demanda}
\end{figure}

Figure NNNNN graphically illustrates how the tool generates a line plot comparing the errors of the selected methods following classification. For this series, Linear Regression exhibits slightly more stable errors compared with ARIMA and Holt, which show greater variability in certain months. Additionally, the three methods tend to underestimate or overestimate at the same points, further confirming the absence of seasonality.

\subsubsection{Parameter Optimization of the Best-Performing Method}
\begin{figure}[H]
    \includegraphics[width=0.8\textwidth]{flowchart_tool.png}
    \caption{Results of parameter optimization}
    \label{fig:demanda}
\end{figure}

Once Linear Regression was identified as the method with the lowest MAPE among the evaluated models for the product, the tool proceeded to verify whether additional adjustments could further improve its performance.
As shown in Figure NNNNN, since the regression method does not require complex hyperparameters, the optimization process confirmed that the initial model already represented the best available alternative. The baseline MAPE was 22.63\%, and after optimization, the same value (22.63\%) was obtained, indicating that the model achieves its best performance without additional adjustments.
These characteristics favor models structured around a trend component, such as Linear Regression, as opposed to more complex methods requiring additional structural patterns to improve fit, such as Holt or ARIMA.

\subsubsection{Forecasting with the Selected Method}

\begin{figure}[H]
    \includegraphics[width=0.8\textwidth]{flowchart_tool.png}
    \caption{Line plot comparing historical data and forecast}
    \label{fig:demanda}
\end{figure}

Using the selected model, the tool generates future forecasts, allowing the user to choose a horizon between 6 and 24 months. Figure 7 presents three main elements of the result:
(i) the historical demand, represented by a blue line showing the observed behavior during 2023–2024;
(ii) the fitted forecast over the historical data, displayed as a red dashed line indicating how the model would have estimated the known portion of the series; and
(iii) the future forecast, shown as a green line projecting a gradual decrease in demand over the next two years.

\begin{figure}[H]
    \includegraphics[width=0.8\textwidth]{flowchart_tool.png}
    \caption{Forecast results displayed in tabular format and download option}
    \label{fig:demanda}
\end{figure}

Figure NNNNN presents each forecasted month along with the expected demand value. For this product, the model suggests a decreasing trajectory, starting at approximately 8,400 units in January 2025 and reaching values close to 3,000 units by the end of 2026.
This information supports decision-making related to inventory planning, raw material purchasing, and production adjustments. Additionally, the tool includes a button labeled “Download forecast (Excel),” allowing users to export the complete forecast in Excel format and directly integrate it into internal analyses without additional processing steps.

\subsubsection{Validation Results}

\begin{figure}[H]
    \includegraphics[width=0.8\textwidth]{flowchart_tool.png}
    \caption{Tool evaluation results}
    \label{fig:demanda}
\end{figure}

\begin{figure}[H]
    \includegraphics[width=0.8\textwidth]{flowchart_tool.png}
    \caption{Tool forecasting results}
    \label{fig:demanda}
\end{figure}

Figures NN and NN present the results obtained from three different types of demand series. Each case includes error comparison plots between the selected models and the corresponding forecast generated using the winning model.

For the series with trend (first column of Figure NN), although both models exhibit very similar error patterns, the blue line (Holt–Winters) shows smoother and more controlled fluctuations compared with the red line (SARIMA), which presents higher peaks in several segments. This is consistent with the performance metrics, where Holt–Winters achieves better MSE, MAD, and MAPE values, indicating more accurate predictions on average. Both the visual error behavior and the numerical indicators therefore support Holt–Winters as the best-performing method for this series.

The forecast generated for the trending series (first column of Figure NNN) using Holt–Winters shows a close fit to the historical demand, evidenced by the near overlap between the actual series (blue) and the fitted forecast (red dashed line). This confirms that the model effectively captures the upward trend and overall behavior of the series. The future projection (green line) continues along a coherent trajectory, maintaining the sustained growth characteristic of the demand pattern. Overall, the forecast is stable and consistent, reinforcing the selection of Holt–Winters.

For the stationary series (second column of Figure NNNN), although the errors of the evaluated methods display very similar behavior over time, the Simple Average method achieves the best metrics, obtaining the lowest MSE, MAD, and MAPE. For a series without trend or seasonality, its simplicity prevents overfitting and directly captures the stable historical behavior. Consequently, both visually and numerically, the Simple Average emerges as the most efficient and accurate method for this case.

The forecast generated using the Simple Average (second column of Figure NNNN) reflects the inherent stability of a stationary series. The red dashed fitted line remains nearly constant around the average demand level, effectively capturing its general behavior. The future projection (green line) continues at the same level, consistent with the nature of the series and avoiding the introduction of artificial patterns. The result is a stable and appropriate forecast, supported by a low adjusted historical MAPE.

For the seasonal series (third column of Figure 9), although both models show similar error patterns over time, the blue line (SARIMA) appears more stable and controlled compared with the red line (Holt–Winters), which exhibits a high initial peak and greater variability. The metrics confirm this observation, as SARIMA achieves better MSE, MAD, and MAPE values, indicating more accurate and consistent predictions on average. Both graphical and numerical evidence therefore identify SARIMA as the most suitable model for this series.

The forecast plot for the seasonal series (third column of Figure NNN) shows that SARIMA effectively captures both the seasonal pattern and the level fluctuations of the historical data. The red dashed fitted line closely follows the observed values, confirming that the model accurately reproduces past dynamics. The future forecast (green line) projects a pattern consistent with the observed seasonality, displaying recurring peaks and troughs without abrupt deviations. Consequently, the model provides a coherent and reliable forecast, supported by an appropriate adjusted historical MAPE for this type of demand.

\subsection{Computational Time}

\begin{table}[H]
\caption{Computational Time of the Tool}
\centering
\begin{tabularx}{\textwidth}{CC}
\toprule
\textbf{Function} & \textbf{Time [s]} \\
\midrule
Metrics and error plot (24 data points) & 1.1 \\
Forecast plot and table (24 data points) & 2.3 \\
Metrics and error plot (48 data points) & 4.2 \\
Forecast plot and table (48 data points) & 7.1 \\
Metrics and error plot (72 data points) & 11.01 \\
Forecast plot and table (72 data points) & 21.06 \\
Metrics and error plot (96 data points) & 16 \\
Forecast plot and table (96 data points) & 26.1 \\
Metrics and error plot (120 data points) & 28.47 \\
Forecast plot and table (120 data points) & 61.32 \\
\bottomrule
\end{tabularx}
\end{table}

Table NNN presents the computational times manually recorded when executing the main functions of the tool for different series sizes. The evaluated functions include the calculation of performance metrics and generation of error plots, as well as the generation of forecast plots and corresponding tables.
As shown in Table 3, computational time increases progressively as the number of observations grows. For small series (24 and 48 data points), processing times remain low, ranging between 1 and 7 seconds, demonstrating efficient performance.
As the dataset size increases (72, 96, and 120 observations), processing times rise accordingly, particularly for functions that combine forecast visualization and tabular outputs. Nevertheless, even in the most demanding evaluated scenario (120 data points), total processing time does not exceed two minutes.
Therefore, although execution time increases with series size, the observed computational cost remains reasonable and does not represent a significant delay for the user.

\subsection{Comparison with Prophet}
To further examine comparative performance, forecast accuracy was analyzed as a function of training sample size. Figure X presents the MAPE values of both Prophet and the proposed tool, ordered by increasing data length, allowing a clear visualization of performance dominance across different information regimes.

The graphical analysis reveals a differentiated performance pattern. For very short series (24 observations), the proposed tool substantially outperforms Prophet, achieving a markedly lower MAPE, indicating greater robustness under data scarcity conditions. In intermediate sample sizes (60–96 observations), the tool maintains competitive performance and frequently achieves lower or comparable error levels. However, as the series length increases beyond 120 observations, Prophet tends to exhibit more stable improvements and, in several cases, achieves lower MAPE values.

These results suggest that the proposed tool demonstrates particular strength in short and moderately sized series, where structural filtering and parsimonious model selection reduce the risk of overfitting. This behavior is consistent with findings from large-scale forecasting competitions, which indicate that simpler or structurally aligned statistical models can outperform more complex approaches when data availability is limited (Makridakis et al., 2020; Makridakis et al., 2022). Conversely, as data volume increases, models designed to capture flexible nonlinear patterns—such as Prophet—may benefit from additional information, leading to improved stability.

Overall, the evidence indicates that the proposed software provides a clear advantage in low-data operational contexts, a common scenario in small and medium-sized inventory systems, while remaining competitive as data availability grows. This reinforces its suitability for real-world demand forecasting environments characterized by heterogeneous series lengths and limited historical information.

%%%%%%%%%%%%%%%%%%%%%%%%%%%%%%%%%%%%%%%%%%
\section{Discussion}

(ELIMINARRRRRR
Authors should discuss the results and how they can be interpreted from the perspective of previous studies and of the working hypotheses. The findings and their implications should be discussed in the broadest context possible. Future research directions may also be highlighted.
ELIMINARRRRR)

The results indicate that the tool effectively selects the most accurate forecasting method according to the structural characteristics of each time series. This finding is consistent with recent forecasting research emphasizing that model performance depends strongly on time-series features such as trend and seasonality, and that selection mechanisms should adapt to structural heterogeneity rather than apply a uniform model across datasets (Montero-Manso et al., 2020; Makridakis et al., 2022).
The ability to compare forecast errors and visualize projected trajectories enhances interpretability and supports operational decision-making processes. Empirical evidence from large-scale forecasting competitions shows that simple statistical models remain highly competitive, particularly in structured and low-noise environments (Makridakis et al., 2020; Petropoulos et al., 2022). This behavior was also observed in several of the analyzed series, where relatively parsimonious approaches outperformed more complex alternatives.
Although computational time increases with longer series, processing remained within acceptable operational limits, supporting the feasibility of real-world implementation. Recent studies highlight that practical forecasting solutions must balance accuracy, interpretability, and computational efficiency to be viable in business environments (Makridakis et al., 2022).
However, the tool presents certain limitations. By relying exclusively on univariate models, it does not incorporate exogenous variables that could potentially improve forecast accuracy in environments influenced by external drivers. Contemporary forecasting literature stresses that multivariate and hybrid approaches can outperform purely univariate specifications when relevant external information is available (Wang et al., 2023).
Furthermore, forecast performance remains dependent on data quality and sample size, which may reduce stability in very short or highly noisy series. Finally, while the automatic selection mechanism is functional, the pool of candidate models is primarily composed of classical statistical approaches and does not include machine learning or deep learning architectures, which have demonstrated competitive performance in complex and large-scale forecasting scenarios (Makridakis et al., 2020; Makridakis et al., 2022).
These limitations provide clear avenues for future development toward a more flexible, hybrid, and scalable forecasting framework.
%%%%%%%%%%%%%%%%%%%%%%%%%%%%%%%%%%%%%%%%%%
\section{Conclusions}

This section is not mandatory, but can be added to the manuscript if the discussion is unusually long or complex.

%%%%%%%%%%%%%%%%%%%%%%%%%%%%%%%%%%%%%%%%%%
\section{Patents}

This section is not mandatory, but may be added if there are patents resulting from the work reported in this manuscript.

%%%%%%%%%%%%%%%%%%%%%%%%%%%%%%%%%%%%%%%%%%
\vspace{6pt} 

%%%%%%%%%%%%%%%%%%%%%%%%%%%%%%%%%%%%%%%%%%
%% optional
%\supplementary{The following supporting information can be downloaded at:  \linksupplementary{s1}, Figure S1: title; Table S1: title; Video S1: title.}

% Only for journal Methods and Protocols:
% If you wish to submit a video article, please do so with any other supplementary material.
% \supplementary{The following supporting information can be downloaded at: \linksupplementary{s1}, Figure S1: title; Table S1: title; Video S1: title. A supporting video article is available at doi: link.}

% Only used for preprtints:
% \supplementary{The following supporting information can be downloaded at the website of this paper posted on \href{https://www.preprints.org/}{Preprints.org}.}

% Only for journal Hardware:
% If you wish to submit a video article, please do so with any other supplementary material.
% \supplementary{The following supporting information can be downloaded at: \linksupplementary{s1}, Figure S1: title; Table S1: title; Video S1: title.\vspace{6pt}\\
%\begin{tabularx}{\textwidth}{lll}
%\toprule
%\textbf{Name} & \textbf{Type} & \textbf{Description} \\
%\midrule
%S1 & Python script (.py) & Script of python source code used in XX \\
%S2 & Text (.txt) & Script of modelling code used to make Figure X \\
%S3 & Text (.txt) & Raw data from experiment X \\
%S4 & Video (.mp4) & Video demonstrating the hardware in use \\
%... & ... & ... \\
%\bottomrule
%\end{tabularx}
%}

%%%%%%%%%%%%%%%%%%%%%%%%%%%%%%%%%%%%%%%%%%
\authorcontributions{For research articles with several authors, a short paragraph specifying their individual contributions must be provided. The following statements should be used ``Conceptualization, X.X. and Y.Y.; methodology, X.X.; software, X.X.; validation, X.X., Y.Y. and Z.Z.; formal analysis, X.X.; investigation, X.X.; resources, X.X.; data curation, X.X.; writing---original draft preparation, X.X.; writing---review and editing, X.X.; visualization, X.X.; supervision, X.X.; project administration, X.X.; funding acquisition, Y.Y. All authors have read and agreed to the published version of the manuscript.'', please turn to the  \href{http://img.mdpi.org/data/contributor-role-instruction.pdf}{CRediT taxonomy} for the term explanation. Authorship must be limited to those who have contributed substantially to the work~reported.}

\funding{Please add: ``This research received no external funding'' or ``This research was funded by NAME OF FUNDER grant number XXX.'' and  and ``The APC was funded by XXX''. Check carefully that the details given are accurate and use the standard spelling of funding agency names at \url{https://search.crossref.org/funding}, any errors may affect your future funding.}

\institutionalreview{In this section, you should add the Institutional Review Board Statement and approval number, if relevant to your study. You might choose to exclude this statement if the study did not require ethical approval. Please note that the Editorial Office might ask you for further information. Please add “The study was conducted in accordance with the Declaration of Helsinki, and approved by the Institutional Review Board (or Ethics Committee) of NAME OF INSTITUTE (protocol code XXX and date of approval).” for studies involving humans. OR “The animal study protocol was approved by the Institutional Review Board (or Ethics Committee) of NAME OF INSTITUTE (protocol code XXX and date of approval).” for studies involving animals. OR “Ethical review and approval were waived for this study due to REASON (please provide a detailed justification).” OR “Not applicable” for studies not involving humans or animals.}

\informedconsent{Any research article describing a study involving humans should contain this statement. Please add ``Informed consent was obtained from all subjects involved in the study.'' OR ``Patient consent was waived due to REASON (please provide a detailed justification).'' OR ``Not applicable'' for studies not involving humans. You might also choose to exclude this statement if the study did not involve humans.

Written informed consent for publication must be obtained from participating patients who can be identified (including by the patients themselves). Please state ``Written informed consent has been obtained from the patient(s) to publish this paper'' if applicable.}

\dataavailability{We encourage all authors of articles published in MDPI journals to share their research data. In this section, please provide details regarding where data supporting reported results can be found, including links to publicly archived datasets analyzed or generated during the study. Where no new data were created, or where data is unavailable due to privacy or ethical restrictions, a statement is still required. Suggested Data Availability Statements are available in section ``MDPI Research Data Policies'' at \url{https://www.mdpi.com/ethics}.} 

%\durcstatement{Current research is limited to the [please insert a specific academic field, e.g., XXX], which is beneficial [share benefits and/or primary use] and does not pose a threat to public health or national security. Authors acknowledge the dual-use potential of the research involving xxx and confirm that all necessary precautions have been taken to prevent potential misuse. As an ethical responsibility, authors strictly adhere to relevant national and international laws about DURC. Authors advocate for responsible deployment, ethical considerations, regulatory compliance, and transparent reporting to mitigate misuse risks and foster beneficial outcomes.}

% Only for journal Nursing Reports
%\publicinvolvement{Please describe how the public (patients, consumers, carers) were involved in the research. Consider reporting against the GRIPP2 (Guidance for Reporting Involvement of Patients and the Public) checklist. If the public were not involved in any aspect of the research add: ``No public involvement in any aspect of this research''.}
%
%% Only for journal Nursing Reports
%\guidelinesstandards{Please add a statement indicating which reporting guideline was used when drafting the report. For example, ``This manuscript was drafted against the XXX (the full name of reporting guidelines and citation) for XXX (type of research) research''. A complete list of reporting guidelines can be accessed via the equator network: \url{https://www.equator-network.org/}.}
%
%% Only for journal Nursing Reports
%\useofartificialintelligence{Please describe in detail any and all uses of artificial intelligence (AI) or AI-assisted tools used in the preparation of the manuscript. This may include, but is not limited to, language translation, language editing and grammar, or generating text. Alternatively, please state that “AI or AI-assisted tools were not used in drafting any aspect of this manuscript”.}

\acknowledgments{In this section you can acknowledge any support given which is not covered by the author contribution or funding sections. This may include administrative and technical support, or donations in kind (e.g., materials used for experiments). Where GenAI has been used for purposes such as generating text, data, or graphics, or for study design, data collection, analysis, or interpretation of data, please add “During the preparation of this manuscript/study, the author(s) used [tool name, version information] for the purposes of [description of use]. The authors have reviewed and edited the output and take full responsibility for the content of this publication.”}

\conflictsofinterest{Declare conflicts of interest or state ``The authors declare no conflicts of interest.'' Authors must identify and declare any personal circumstances or interest that may be perceived as inappropriately influencing the representation or interpretation of reported research results. Any role of the funders in the design of the study; in the collection, analyses or interpretation of data; in the writing of the manuscript; or in the decision to publish the results must be declared in this section. If there is no role, please state ``The funders had no role in the design of the study; in the collection, analyses, or interpretation of data; in the writing of the manuscript; or in the decision to publish the results''.} 

%%%%%%%%%%%%%%%%%%%%%%%%%%%%%%%%%%%%%%%%%%
%% Optional

%% Only for journal Encyclopedia
%\entrylink{The Link to this entry published on the encyclopedia platform.}

\abbreviations{Abbreviations}{
The following abbreviations are used in this manuscript:
\\

\noindent 
\begin{tabular}{@{}ll}
MDPI & Multidisciplinary Digital Publishing Institute\\
DOAJ & Directory of open access journals\\
TLA & Three letter acronym\\
LD & Linear dichroism
\end{tabular}
}

%%%%%%%%%%%%%%%%%%%%%%%%%%%%%%%%%%%%%%%%%%
%% Optional
\appendixtitles{no} % Leave argument "no" if all appendix headings stay EMPTY (then no dot is printed after "Appendix A"). If the appendix sections contain a heading then change the argument to "yes".
\appendixstart
\appendix
\section[\appendixname~\thesection]{}
\subsection[\appendixname~\thesubsection]{}
The appendix is an optional section that can contain details and data supplemental to the main text---for example, explanations of experimental details that would disrupt the flow of the main text but nonetheless remain crucial to understanding and reproducing the research shown; figures of replicates for experiments of which representative data are shown in the main text can be added here if brief, or as Supplementary Data. Mathematical proofs of results not central to the paper can be added as an appendix.

\begin{table}[H] 
\caption{This is a table caption.\label{tab5}}
%\newcolumntype{C}{>{\centering\arraybackslash}X}
\begin{tabularx}{\textwidth}{CCC}
\toprule
\textbf{Title 1}	& \textbf{Title 2}	& \textbf{Title 3}\\
\midrule
Entry 1		& Data			& Data\\
Entry 2		& Data			& Data\\
\bottomrule
\end{tabularx}
\end{table}

\section[\appendixname~\thesection]{}
All appendix sections must be cited in the main text. In the appendices, Figures, Tables, etc. should be labeled, starting with ``A''---e.g., Figure A1, Figure A2, etc.

%%%%%%%%%%%%%%%%%%%%%%%%%%%%%%%%%%%%%%%%%%
%\isPreprints{}{% This command is only used for ``preprints''.
\begin{adjustwidth}{-\extralength}{0cm}
%} % If the paper is ``preprints'', please uncomment this parenthesis.
%\printendnotes[custom] % Un-comment to print a list of endnotes

\reftitle{References}

% Please provide the correct journal abbreviation (e.g. according to the “List of Title Word Abbreviations” http://www.issn.org/services/online-services/access-to-the-ltwa/).
% Citations and References in Supplementary files are permitted provided that they also appear in the reference list here. 

%=====================================
% References, variant A: external bibliography
%=====================================
% \bibliography{your_external_BibTeX_file}

%=====================================
% References, variant B: internal bibliography
%=====================================

% ACS format
\begin{thebibliography}{999}
% Reference 1
\bibitem{ref-journal}
Author~1, T. The title of the cited article. {\em Journal Abbreviation} {\bf 2008}, {\em 10}, 142--149.
% Reference 2
\bibitem{ref-book1}
Author~2, L. The title of the cited contribution. In {\em The Book Title}; Editor 1, F., Editor 2, A., Eds.; Publishing House: City, Country, 2007; pp. 32--58.
% Reference 3
\bibitem{ref-book2}
Author 1, A.; Author 2, B. \textit{Book Title}, 3rd ed.; Publisher: Publisher Location, Country, 2008; pp. 154--196.
% Reference 4
\bibitem{ref-unpublish}
Author 1, A.B.; Author 2, C. Title of Unpublished Work. \textit{Abbreviated Journal Name} year, \textit{phrase indicating stage of publication (submitted; accepted; in press)}.
% Reference 5
\bibitem{ref-url}
Title of Site. Available online: URL (accessed on Day Month Year).
% Reference 6
\bibitem{ref-proceeding}
Author 1, A.B.; Author 2, C.D.; Author 3, E.F. Title of presentation. In Proceedings of the Name of the Conference, Location of Conference, Country, Date of Conference (Day Month Year); Abstract Number (optional), Pagination (optional).
% Reference 7
\bibitem{ref-thesis}
Author 1, A.B. Title of Thesis. Level of Thesis, Degree-Granting University, Location of University, Date of Completion.
\end{thebibliography}

% If authors have biography, please use the format below
%\section*{Short Biography of Authors}
%\bio
%{\raisebox{-0.35cm}{\includegraphics[width=3.5cm,height=5.3cm,clip,keepaspectratio]{Definitions/author1.pdf}}}
%{\textbf{Firstname Lastname} Biography of first author}
%
%\bio
%{\raisebox{-0.35cm}{\includegraphics[width=3.5cm,height=5.3cm,clip,keepaspectratio]{Definitions/author2.jpg}}}
%{\textbf{Firstname Lastname} Biography of second author}

% For the MDPI journals use author-date citation, please follow the formatting guidelines on http://www.mdpi.com/authors/references
% To cite two works by the same author: \citeauthor{ref-journal-1a} (\citeyear{ref-journal-1a}, \citeyear{ref-journal-1b}). This produces: Whittaker (1967, 1975)
% To cite two works by the same author with specific pages: \citeauthor{ref-journal-3a} (\citeyear{ref-journal-3a}, p. 328; \citeyear{ref-journal-3b}, p.475). This produces: Wong (1999, p. 328; 2000, p. 475)

%%%%%%%%%%%%%%%%%%%%%%%%%%%%%%%%%%%%%%%%%%
%% for journal Sci
%\reviewreports{\\
%Reviewer 1 comments and authors’ response\\
%Reviewer 2 comments and authors’ response\\
%Reviewer 3 comments and authors’ response
%}
%%%%%%%%%%%%%%%%%%%%%%%%%%%%%%%%%%%%%%%%%%
\PublishersNote{}
%\isPreprints{}{% This command is only used for ``preprints''.
\end{adjustwidth}
%} % If the paper is ``preprints'', please uncomment this parenthesis.
\end{document}

